%% ---------------------------------------------------------------------
\section{Build instructions}
\label{sec:build}

\noindent\fcolorbox{red!60!black}{red!5}{\parbox{\dimexpr\textwidth-2\fboxsep-2\fboxrule}{\small
\textbf{Safety.} Remove all propellers during assembly and bench testing. Never power
the motors with the companion computer connected for the first time. Check the RC
override and failsafes before every autonomous flight. Lithium-polymer batteries
must be charged and stored per the manufacturer's instructions.}}

\subsection{Mechanical integration}
\label{sec:build-mech}

\begin{enumerate}[leftmargin=*,itemsep=2pt]
  \item Mount the flight controller on vibration isolation at the airframe's centre
    of gravity, arrow forward, as for any ArduPilot build. Confirm normal manual
    flight \emph{before} adding the companion.
  \item Mount the Raspberry Pi where its USB ports are reachable and it receives
    airflow. A Pi under sustained load needs a heatsink.
  \item Mount the YDLIDAR G2 with an unobstructed horizontal field of view. Propeller
    arms and landing gear in the scan plane appear as permanent obstacles. Returns
    closer than \SI{0.10}{\metre} (costmap) and \SI{0.12}{\metre} (streamer) are
    rejected, but larger airframe parts must be kept out of the plane or masked.
    Record the mount offset (\cref{sec:build-calib}).
  \item Mount the downward optical-flow/range module with a clear view of the floor,
    away from propeller downwash and legs, oriented per ArduPilot's
    \code{FLOW_ORIENT_YAW} convention.
  \item Mount the camera facing forward and record its offset.
  \item Route the USB cables with strain relief. A connector that is pulled loose in
    flight removes that sensor; the supervisor restarts the stream when the device
    returns.
\end{enumerate}

\subsection{Electrical integration}

Follow \cref{fig:wiring}. Connect the autopilot to the Pi by USB (simplest and
recommended). Connect the LiDAR adapter and camera to the Pi's other USB ports. Supply
the Pi from a dedicated \SI{5}{\volt} regulator rated at least \SI{3}{\ampere}. The
flow/range module and RC receiver connect to the autopilot as documented by
ArduPilot.

\subsection{Companion computer provisioning}

On the Pi, which needs only the \code{raspberry_pi/} directory:
\begin{lstlisting}[language=bash]
# from the laptop: copy code (repeat after every change)
rsync -av --exclude '*.pt' --exclude '__pycache__' \
  ~/gps-denied-vtol/raspberry_pi/ <user>@<pi_ip>:~/raspberry_pi/

# on the Pi
cd ~/raspberry_pi
sudo apt install -y python3-opencv python3-numpy
pip install -r requirements.txt --break-system-packages
sudo ./setup_pi.sh            # udev rules, dialout, chrony, iw, sudoers entries
# log out and in once, then verify:
ls -la /dev/pixhawk /dev/camera /dev/ttyUSB0
python3 test_connection.py    # must end with [PASSED]
\end{lstlisting}
Other autopilot or camera models are supported by adding their USB vendor/product IDs
to \code{99-pixhawk.rules} or \code{98-camera.rules}.

\subsection{Ground computer provisioning}

Requirements: Ubuntu 22.04, ROS~2 Humble and an NVIDIA GPU.
\begin{lstlisting}[language=bash]
sudo apt install -y python3-colcon-common-extensions python3-rosdep
sudo rosdep init; rosdep update
git clone https://github.com/DEEPAKSINGLA-byte/gps-denied-vtol.git ~/gps-denied-vtol
cd ~/gps-denied-vtol/ros2_ws && source /opt/ros/humble/setup.bash
rosdep install --from-paths src --ignore-src -r -y
colcon build --symlink-install
pip install -U ultralytics                 # YOLO + SAM

# disable Wi-Fi power saving and serve time to the Pi
printf '[connection]\nwifi.powersave = 2\n' | \
  sudo tee /etc/NetworkManager/conf.d/wifi-powersave-off.conf
sudo systemctl restart NetworkManager
sudo ~/gps-denied-vtol/scripts/setup_laptop_timeserver.sh
\end{lstlisting}
Model weights (\code{yolov8n-seg.pt}, \code{sam2.1_t.pt}) are resolved from the
workspace root, so they are not re-downloaded at every launch. A laptop connected to
the drone's hotspot is usually offline, and the segmentation node would otherwise
fail. Connect Pi and laptop directly (e.g.\ a \SI{5}{\giga\hertz} laptop hotspot).
Relaying through a phone hotspot sends each packet over the air twice.

\subsection{Autopilot configuration}

\begin{table}[H]
\centering\small
\caption{ArduPilot parameters for indoor GPS-denied operation (\code{docs/SETUP.md}).}
\label{tab:apparams}
\begin{tabularx}{\textwidth}{l l L}
\toprule
\textbf{Parameter} & \textbf{Value} & \textbf{Reason} \\
\midrule
\code{FLTMODEx} & \code{GUIDED} (4) on one switch position & The laptop can steer only in GUIDED \\
\code{EK3_SRC1_POSXY} & 0 (none) or 4 (optical flow) & No GNSS indoors \\
\code{EK3_SRC1_POSZ} & 1 (baro) or 2 (rangefinder) & Height source \\
\code{AHRS_GPS_USE} & 0 & Do not use GPS \\
\code{ARMING_CHECK} & 1048558 & Skips GPS/compass checks that otherwise block arming (\code{PreArm: Need 3D Fix}) \\
RC-loss, battery failsafes & per vehicle & Must land/hold without GNSS \\
\bottomrule
\end{tabularx}
\end{table}

\subsection{Sensor extrinsic calibration}
\label{sec:build-calib}

The mount transforms in \code{navigation_real.launch.py} are placeholders:
LiDAR $(0,0,-0.15)$\,m and camera $(0.10,0,0.05)$\,m, both with zero rotation.
\begin{enumerate}[leftmargin=*,itemsep=1pt]
  \item Measure the translation from the flight controller's centre to the LiDAR's
    optical centre and to the camera, in FLU axes.
  \item Measure the LiDAR's yaw relative to the airframe's forward axis. The receiver
    also exposes \code{angle_offset_deg} and \code{clockwise_input}.
  \item Edit the static transform arguments, rebuild (\code{colcon build}), relaunch.
  \item Validate. With the vehicle carried slowly by hand past walls, the scan must
    stay registered to the costmap. \topic{/metrics/costmap/local/scan_agreement}
    should remain close to 1, and a value well below 1 indicates a wrong mount yaw or
    timing.
\end{enumerate}
Optical-flow sign errors are corrected on the autopilot (\code{FLOW_ORIENT_YAW}), or
for a pure $x$ flip with the bridge parameter \code{optical_flow_scale}$=-1$.
